% Options for packages loaded elsewhere
\PassOptionsToPackage{unicode}{hyperref}
\PassOptionsToPackage{hyphens}{url}
\documentclass[
]{article}
\usepackage{xcolor}
\usepackage[margin=1in]{geometry}
\usepackage{amsmath,amssymb}
\setcounter{secnumdepth}{-\maxdimen} % remove section numbering
\usepackage{iftex}
\ifPDFTeX
  \usepackage[T1]{fontenc}
  \usepackage[utf8]{inputenc}
  \usepackage{textcomp} % provide euro and other symbols
\else % if luatex or xetex
  \usepackage{unicode-math} % this also loads fontspec
  \defaultfontfeatures{Scale=MatchLowercase}
  \defaultfontfeatures[\rmfamily]{Ligatures=TeX,Scale=1}
\fi
\usepackage{lmodern}
\ifPDFTeX\else
  % xetex/luatex font selection
\fi
% Use upquote if available, for straight quotes in verbatim environments
\IfFileExists{upquote.sty}{\usepackage{upquote}}{}
\IfFileExists{microtype.sty}{% use microtype if available
  \usepackage[]{microtype}
  \UseMicrotypeSet[protrusion]{basicmath} % disable protrusion for tt fonts
}{}
\makeatletter
\@ifundefined{KOMAClassName}{% if non-KOMA class
  \IfFileExists{parskip.sty}{%
    \usepackage{parskip}
  }{% else
    \setlength{\parindent}{0pt}
    \setlength{\parskip}{6pt plus 2pt minus 1pt}}
}{% if KOMA class
  \KOMAoptions{parskip=half}}
\makeatother
\usepackage{color}
\usepackage{fancyvrb}
\newcommand{\VerbBar}{|}
\newcommand{\VERB}{\Verb[commandchars=\\\{\}]}
\DefineVerbatimEnvironment{Highlighting}{Verbatim}{commandchars=\\\{\}}
% Add ',fontsize=\small' for more characters per line
\usepackage{framed}
\definecolor{shadecolor}{RGB}{248,248,248}
\newenvironment{Shaded}{\begin{snugshade}}{\end{snugshade}}
\newcommand{\AlertTok}[1]{\textcolor[rgb]{0.94,0.16,0.16}{#1}}
\newcommand{\AnnotationTok}[1]{\textcolor[rgb]{0.56,0.35,0.01}{\textbf{\textit{#1}}}}
\newcommand{\AttributeTok}[1]{\textcolor[rgb]{0.13,0.29,0.53}{#1}}
\newcommand{\BaseNTok}[1]{\textcolor[rgb]{0.00,0.00,0.81}{#1}}
\newcommand{\BuiltInTok}[1]{#1}
\newcommand{\CharTok}[1]{\textcolor[rgb]{0.31,0.60,0.02}{#1}}
\newcommand{\CommentTok}[1]{\textcolor[rgb]{0.56,0.35,0.01}{\textit{#1}}}
\newcommand{\CommentVarTok}[1]{\textcolor[rgb]{0.56,0.35,0.01}{\textbf{\textit{#1}}}}
\newcommand{\ConstantTok}[1]{\textcolor[rgb]{0.56,0.35,0.01}{#1}}
\newcommand{\ControlFlowTok}[1]{\textcolor[rgb]{0.13,0.29,0.53}{\textbf{#1}}}
\newcommand{\DataTypeTok}[1]{\textcolor[rgb]{0.13,0.29,0.53}{#1}}
\newcommand{\DecValTok}[1]{\textcolor[rgb]{0.00,0.00,0.81}{#1}}
\newcommand{\DocumentationTok}[1]{\textcolor[rgb]{0.56,0.35,0.01}{\textbf{\textit{#1}}}}
\newcommand{\ErrorTok}[1]{\textcolor[rgb]{0.64,0.00,0.00}{\textbf{#1}}}
\newcommand{\ExtensionTok}[1]{#1}
\newcommand{\FloatTok}[1]{\textcolor[rgb]{0.00,0.00,0.81}{#1}}
\newcommand{\FunctionTok}[1]{\textcolor[rgb]{0.13,0.29,0.53}{\textbf{#1}}}
\newcommand{\ImportTok}[1]{#1}
\newcommand{\InformationTok}[1]{\textcolor[rgb]{0.56,0.35,0.01}{\textbf{\textit{#1}}}}
\newcommand{\KeywordTok}[1]{\textcolor[rgb]{0.13,0.29,0.53}{\textbf{#1}}}
\newcommand{\NormalTok}[1]{#1}
\newcommand{\OperatorTok}[1]{\textcolor[rgb]{0.81,0.36,0.00}{\textbf{#1}}}
\newcommand{\OtherTok}[1]{\textcolor[rgb]{0.56,0.35,0.01}{#1}}
\newcommand{\PreprocessorTok}[1]{\textcolor[rgb]{0.56,0.35,0.01}{\textit{#1}}}
\newcommand{\RegionMarkerTok}[1]{#1}
\newcommand{\SpecialCharTok}[1]{\textcolor[rgb]{0.81,0.36,0.00}{\textbf{#1}}}
\newcommand{\SpecialStringTok}[1]{\textcolor[rgb]{0.31,0.60,0.02}{#1}}
\newcommand{\StringTok}[1]{\textcolor[rgb]{0.31,0.60,0.02}{#1}}
\newcommand{\VariableTok}[1]{\textcolor[rgb]{0.00,0.00,0.00}{#1}}
\newcommand{\VerbatimStringTok}[1]{\textcolor[rgb]{0.31,0.60,0.02}{#1}}
\newcommand{\WarningTok}[1]{\textcolor[rgb]{0.56,0.35,0.01}{\textbf{\textit{#1}}}}
\usepackage{longtable,booktabs,array}
\newcounter{none} % for unnumbered tables
\usepackage{calc} % for calculating minipage widths
% Correct order of tables after \paragraph or \subparagraph
\usepackage{etoolbox}
\makeatletter
\patchcmd\longtable{\par}{\if@noskipsec\mbox{}\fi\par}{}{}
\makeatother
% Allow footnotes in longtable head/foot
\IfFileExists{footnotehyper.sty}{\usepackage{footnotehyper}}{\usepackage{footnote}}
\makesavenoteenv{longtable}
\usepackage{graphicx}
\makeatletter
\newsavebox\pandoc@box
\newcommand*\pandocbounded[1]{% scales image to fit in text height/width
  \sbox\pandoc@box{#1}%
  \Gscale@div\@tempa{\textheight}{\dimexpr\ht\pandoc@box+\dp\pandoc@box\relax}%
  \Gscale@div\@tempb{\linewidth}{\wd\pandoc@box}%
  \ifdim\@tempb\p@<\@tempa\p@\let\@tempa\@tempb\fi% select the smaller of both
  \ifdim\@tempa\p@<\p@\scalebox{\@tempa}{\usebox\pandoc@box}%
  \else\usebox{\pandoc@box}%
  \fi%
}
% Set default figure placement to htbp
\def\fps@figure{htbp}
\makeatother
\setlength{\emergencystretch}{3em} % prevent overfull lines
\providecommand{\tightlist}{%
  \setlength{\itemsep}{0pt}\setlength{\parskip}{0pt}}
\usepackage{bookmark}
\IfFileExists{xurl.sty}{\usepackage{xurl}}{} % add URL line breaks if available
\urlstyle{same}
\hypersetup{
  pdftitle={Análisis Multidimensional de Indicadores Educativos},
  pdfauthor={Reporte Estadístico},
  hidelinks,
  pdfcreator={LaTeX via pandoc}}

\title{Análisis Multidimensional de Indicadores Educativos}
\author{Reporte Estadístico}
\date{2026-08-28}

\begin{document}
\maketitle

{
\setcounter{tocdepth}{3}
\tableofcontents
}
\subsection{Resumen del Dataset}\label{resumen-del-dataset}

Este reporte presenta una evaluación completa de múltiples variables e
indicadores de cobertura, aprobación, reprobación, repitencia y
deserción del sector educativo. El análisis discrimina adecuadamente las
variables categóricas de las numéricas, normaliza los datos continuos y
explora las distribuciones, valores atípicos y correlaciones
estadísticas de interés.

\begin{center}\rule{0.5\linewidth}{0.5pt}\end{center}

\subsection{1. Preparación del Entorno y Carga de
Datos}\label{preparaciuxf3n-del-entorno-y-carga-de-datos}

En esta etapa inicial, se importan los datos aplicando una distinción
estricta de los separadores de archivo (comas para dividir columnas) y
los delimitadores de decimales (puntos). Luego, con base en el
entendimiento funcional (diccionario de datos), reclasificamos las
variables geográficas que originalmente son numéricas (Códigos DANE)
como categóricas (factores).

\begin{Shaded}
\begin{Highlighting}[]
\CommentTok{\# Cargar librerías necesarias}
\FunctionTok{library}\NormalTok{(ggplot2)}
\FunctionTok{library}\NormalTok{(dplyr)}
\FunctionTok{library}\NormalTok{(tidyr)}
\FunctionTok{library}\NormalTok{(reshape2)}
\FunctionTok{library}\NormalTok{(knitr)}

\CommentTok{\# Crear carpeta para guardar todos los resultados}
\FunctionTok{dir.create}\NormalTok{(}\StringTok{"resultados taller 1"}\NormalTok{, }\AttributeTok{showWarnings =} \ConstantTok{FALSE}\NormalTok{)}

\CommentTok{\# 1. Carga de datos con parámetros explícitos (sep para comas de archivo, dec para decimales)}
\NormalTok{datos }\OtherTok{\textless{}{-}} \FunctionTok{read.csv}\NormalTok{(}\StringTok{"data/EDUCACION\_LIMPIO.csv"}\NormalTok{, }\AttributeTok{sep =} \StringTok{","}\NormalTok{, }\AttributeTok{dec =} \StringTok{"."}\NormalTok{)}

\CommentTok{\# 2. Análisis y Tipificación según "Diccionario de Datos":}
\CommentTok{\# Convertimos las entidades geográficas a factores (categóricas) ya que no son magnitudes numéricas reales.}
\NormalTok{datos}\SpecialCharTok{$}\NormalTok{CÓDIGO\_MUNICIPIO }\OtherTok{\textless{}{-}} \FunctionTok{as.factor}\NormalTok{(datos}\SpecialCharTok{$}\NormalTok{CÓDIGO\_MUNICIPIO)}
\NormalTok{datos}\SpecialCharTok{$}\NormalTok{CÓDIGO\_DEPARTAMENTO }\OtherTok{\textless{}{-}} \FunctionTok{as.factor}\NormalTok{(datos}\SpecialCharTok{$}\NormalTok{CÓDIGO\_DEPARTAMENTO)}

\CommentTok{\# Filtramos estrictamente las variables numéricas (cuantitativas reales) para el análisis estadístico}
\NormalTok{variables\_numericas }\OtherTok{\textless{}{-}}\NormalTok{ datos }\SpecialCharTok{\%\textgreater{}\%} \FunctionTok{select}\NormalTok{(}\FunctionTok{where}\NormalTok{(is.numeric))}
\end{Highlighting}
\end{Shaded}

\begin{center}\rule{0.5\linewidth}{0.5pt}\end{center}

\subsection{2. Tabla de Resumen Estadístico (Variables
Clave)}\label{tabla-de-resumen-estaduxedstico-variables-clave}

La siguiente tabla consolida las estadísticas descriptivas para todas
las variables cuantitativas del modelo.

\begin{Shaded}
\begin{Highlighting}[]
\CommentTok{\# Cálculo de estadísticas clave}
\NormalTok{resumen }\OtherTok{\textless{}{-}} \FunctionTok{data.frame}\NormalTok{(}
  \AttributeTok{Variable =} \FunctionTok{names}\NormalTok{(variables\_numericas),}
  \AttributeTok{Media =} \FunctionTok{sapply}\NormalTok{(variables\_numericas, mean, }\AttributeTok{na.rm =} \ConstantTok{TRUE}\NormalTok{),}
  \AttributeTok{Desv\_Est =} \FunctionTok{sapply}\NormalTok{(variables\_numericas, sd, }\AttributeTok{na.rm =} \ConstantTok{TRUE}\NormalTok{),}
\NormalTok{  Mínimo }\OtherTok{=} \FunctionTok{sapply}\NormalTok{(variables\_numericas, min, }\AttributeTok{na.rm =} \ConstantTok{TRUE}\NormalTok{),}
  \AttributeTok{Mediana =} \FunctionTok{sapply}\NormalTok{(variables\_numericas, median, }\AttributeTok{na.rm =} \ConstantTok{TRUE}\NormalTok{),}
\NormalTok{  Máximo }\OtherTok{=} \FunctionTok{sapply}\NormalTok{(variables\_numericas, max, }\AttributeTok{na.rm =} \ConstantTok{TRUE}\NormalTok{)}
\NormalTok{)}

\CommentTok{\# Exportar los resultados a CSV}
\FunctionTok{write.csv}\NormalTok{(resumen, }\StringTok{"resultados taller 1/Resumen\_Estadistico.csv"}\NormalTok{, }\AttributeTok{row.names =} \ConstantTok{FALSE}\NormalTok{)}

\CommentTok{\# Renderización de tabla en HTML}
\FunctionTok{kable}\NormalTok{(resumen, }\AttributeTok{digits =} \DecValTok{3}\NormalTok{, }\AttributeTok{caption =} \StringTok{"Resumen Estadístico de las Variables Numéricas"}\NormalTok{, }\AttributeTok{row.names =} \ConstantTok{FALSE}\NormalTok{)}
\end{Highlighting}
\end{Shaded}

\begin{longtable}[]{@{}
  >{\raggedright\arraybackslash}p{(\linewidth - 10\tabcolsep) * \real{0.3803}}
  >{\raggedleft\arraybackslash}p{(\linewidth - 10\tabcolsep) * \real{0.1268}}
  >{\raggedleft\arraybackslash}p{(\linewidth - 10\tabcolsep) * \real{0.1268}}
  >{\raggedleft\arraybackslash}p{(\linewidth - 10\tabcolsep) * \real{0.0986}}
  >{\raggedleft\arraybackslash}p{(\linewidth - 10\tabcolsep) * \real{0.1268}}
  >{\raggedleft\arraybackslash}p{(\linewidth - 10\tabcolsep) * \real{0.1408}}@{}}
\caption{Resumen Estadístico de las Variables Numéricas}\tabularnewline
\toprule\noalign{}
\begin{minipage}[b]{\linewidth}\raggedright
Variable
\end{minipage} & \begin{minipage}[b]{\linewidth}\raggedleft
Media
\end{minipage} & \begin{minipage}[b]{\linewidth}\raggedleft
Desv\_Est
\end{minipage} & \begin{minipage}[b]{\linewidth}\raggedleft
Mínimo
\end{minipage} & \begin{minipage}[b]{\linewidth}\raggedleft
Mediana
\end{minipage} & \begin{minipage}[b]{\linewidth}\raggedleft
Máximo
\end{minipage} \\
\midrule\noalign{}
\endfirsthead
\toprule\noalign{}
\begin{minipage}[b]{\linewidth}\raggedright
Variable
\end{minipage} & \begin{minipage}[b]{\linewidth}\raggedleft
Media
\end{minipage} & \begin{minipage}[b]{\linewidth}\raggedleft
Desv\_Est
\end{minipage} & \begin{minipage}[b]{\linewidth}\raggedleft
Mínimo
\end{minipage} & \begin{minipage}[b]{\linewidth}\raggedleft
Mediana
\end{minipage} & \begin{minipage}[b]{\linewidth}\raggedleft
Máximo
\end{minipage} \\
\midrule\noalign{}
\endhead
\bottomrule\noalign{}
\endlastfoot
POBLACIÓN\_5\_16 & 3143.451 & 2819.803 & 1.005 & 2220.000 & 14013.000 \\
TASA\_MATRICULACIÓN\_5\_16 & 0.832 & 0.127 & 0.454 & 0.837 & 1.235 \\
COBERTURA\_NETA & 0.841 & 0.115 & 0.513 & 0.851 & 1.170 \\
COBERTURA\_NETA\_TRANSICIÓN & 0.575 & 0.134 & 0.206 & 0.575 & 0.946 \\
COBERTURA\_NETA\_PRIMARIA & 0.810 & 0.120 & 0.479 & 0.813 & 1.156 \\
COBERTURA\_NETA\_SECUNDARIA & 0.715 & 0.130 & 0.342 & 0.715 & 1.076 \\
COBERTURA\_NETA\_MEDIA & 0.431 & 0.121 & 0.076 & 0.435 & 0.762 \\
COBERTURA\_BRUTA & 0.947 & 0.141 & 0.552 & 0.953 & 1.356 \\
COBERTURA\_BRUTA\_TRANSICIÓN & 0.822 & 0.177 & 0.348 & 0.819 & 1.348 \\
COBERTURA\_BRUTA\_PRIMARIA & 1.001 & 0.176 & 0.534 & 1.000 & 1.478 \\
COBERTURA\_BRUTA\_SECUNDARIA & 1.006 & 0.181 & 0.516 & 1.006 & 1.497 \\
COBERTURA\_BRUTA\_MEDIA & 0.758 & 0.187 & 0.237 & 0.758 & 1.280 \\
DESERCIÓN & 0.030 & 0.016 & 0.000 & 0.028 & 0.083 \\
DESERCIÓN\_TRANSICIÓN & 0.029 & 0.021 & 0.000 & 0.026 & 0.095 \\
DESERCIÓN\_PRIMARIA & 0.023 & 0.015 & 0.000 & 0.020 & 0.070 \\
DESERCIÓN\_SECUNDARIA & 0.041 & 0.024 & 0.000 & 0.037 & 0.114 \\
DESERCIÓN\_MEDIA & 0.028 & 0.019 & 0.000 & 0.025 & 0.084 \\
APROBACIÓN & 0.934 & 0.037 & 0.809 & 0.936 & 1.000 \\
APROBACIÓN\_TRANSICIÓN & 0.000 & 0.000 & 0.000 & 0.000 & 0.000 \\
APROBACIÓN\_PRIMARIA & 0.948 & 0.034 & 0.832 & 0.953 & 1.000 \\
APROBACIÓN\_SECUNDARIA & 0.905 & 0.059 & 0.698 & 0.911 & 1.000 \\
APROBACIÓN\_MEDIA & 0.941 & 0.037 & 0.818 & 0.946 & 1.000 \\
REPROBACIÓN & 0.036 & 0.031 & 0.000 & 0.034 & 0.150 \\
REPROBACIÓN\_TRANSICIÓN & 0.000 & 0.000 & 0.000 & 0.000 & 0.000 \\
REPROBACIÓN\_PRIMARIA & 0.029 & 0.028 & 0.000 & 0.023 & 0.126 \\
REPROBACIÓN\_SECUNDARIA & 0.055 & 0.051 & 0.000 & 0.048 & 0.239 \\
REPROBACIÓN\_MEDIA & 0.031 & 0.030 & 0.000 & 0.023 & 0.125 \\
REPITENCIA & 0.024 & 0.022 & 0.000 & 0.018 & 0.101 \\
REPITENCIA\_TRANSICIÓN & 0.003 & 0.006 & 0.000 & 0.000 & 0.024 \\
REPITENCIA\_PRIMARIA & 0.023 & 0.023 & 0.000 & 0.016 & 0.099 \\
REPITENCIA\_SECUNDARIA & 0.034 & 0.034 & 0.000 & 0.023 & 0.141 \\
REPITENCIA\_MEDIA & 0.012 & 0.013 & 0.000 & 0.007 & 0.053 \\
\end{longtable}

\begin{center}\rule{0.5\linewidth}{0.5pt}\end{center}

\subsection{3. Histogramas para las Variables
Numéricas}\label{histogramas-para-las-variables-numuxe9ricas}

Se presentan las distribuciones individuales para evaluar la simetría,
sesgos y concentración en cada indicador.

\begin{Shaded}
\begin{Highlighting}[]
\CommentTok{\# Transformar datos a formato largo para facilitar la graficación con facetas}
\NormalTok{datos\_long }\OtherTok{\textless{}{-}}\NormalTok{ variables\_numericas }\SpecialCharTok{\%\textgreater{}\%} 
  \FunctionTok{pivot\_longer}\NormalTok{(}\AttributeTok{cols =} \FunctionTok{everything}\NormalTok{(), }\AttributeTok{names\_to =} \StringTok{"Variable"}\NormalTok{, }\AttributeTok{values\_to =} \StringTok{"Valor"}\NormalTok{)}

\CommentTok{\# Graficar la cuadrícula de histogramas}
\NormalTok{p\_hist }\OtherTok{\textless{}{-}} \FunctionTok{ggplot}\NormalTok{(datos\_long, }\FunctionTok{aes}\NormalTok{(}\AttributeTok{x =}\NormalTok{ Valor)) }\SpecialCharTok{+}
  \FunctionTok{geom\_histogram}\NormalTok{(}\AttributeTok{bins =} \DecValTok{30}\NormalTok{, }\AttributeTok{fill =} \StringTok{"\#3498db"}\NormalTok{, }\AttributeTok{color =} \StringTok{"black"}\NormalTok{, }\AttributeTok{alpha =} \FloatTok{0.8}\NormalTok{) }\SpecialCharTok{+}
  \FunctionTok{facet\_wrap}\NormalTok{(}\SpecialCharTok{\textasciitilde{}}\NormalTok{ Variable, }\AttributeTok{scales =} \StringTok{"free"}\NormalTok{, }\AttributeTok{ncol =} \DecValTok{4}\NormalTok{) }\SpecialCharTok{+}
  \FunctionTok{theme\_minimal}\NormalTok{() }\SpecialCharTok{+}
  \FunctionTok{labs}\NormalTok{(}
    \AttributeTok{title =} \StringTok{"Distribución de Variables Educativas (Histogramas)"}\NormalTok{,}
    \AttributeTok{x =} \StringTok{"Valor del Indicador"}\NormalTok{,}
    \AttributeTok{y =} \StringTok{"Frecuencia"}
\NormalTok{  ) }\SpecialCharTok{+}
  \FunctionTok{theme}\NormalTok{(}\AttributeTok{strip.text =} \FunctionTok{element\_text}\NormalTok{(}\AttributeTok{size =} \DecValTok{8}\NormalTok{, }\AttributeTok{face =} \StringTok{"bold"}\NormalTok{))}

\CommentTok{\# Mostrar y guardar gráfico}
\FunctionTok{print}\NormalTok{(p\_hist)}
\end{Highlighting}
\end{Shaded}

\begin{center}\includegraphics{Taller1_files/figure-latex/histogramas-1} \end{center}

\begin{Shaded}
\begin{Highlighting}[]
\FunctionTok{ggsave}\NormalTok{(}\StringTok{"resultados taller 1/Histogramas.png"}\NormalTok{, }\AttributeTok{plot =}\NormalTok{ p\_hist, }\AttributeTok{width =} \DecValTok{12}\NormalTok{, }\AttributeTok{height =} \DecValTok{10}\NormalTok{)}
\end{Highlighting}
\end{Shaded}

\begin{center}\rule{0.5\linewidth}{0.5pt}\end{center}

\subsection{4. Normalización y Diagramas de Caja
(Boxplots)}\label{normalizaciuxf3n-y-diagramas-de-caja-boxplots}

Para evaluar los indicadores en una escala comparable e identificar la
presencia de valores atípicos (outliers), se aplica una normalización
Min-Max (rango de {[}0, 1{]}) y se consolidan en diagramas de caja.

\begin{Shaded}
\begin{Highlighting}[]
\CommentTok{\# Función de normalización Min{-}Max}
\NormalTok{normalizar }\OtherTok{\textless{}{-}} \ControlFlowTok{function}\NormalTok{(x) \{}
  \FunctionTok{return}\NormalTok{ ((x }\SpecialCharTok{{-}} \FunctionTok{min}\NormalTok{(x, }\AttributeTok{na.rm =} \ConstantTok{TRUE}\NormalTok{)) }\SpecialCharTok{/}\NormalTok{ (}\FunctionTok{max}\NormalTok{(x, }\AttributeTok{na.rm =} \ConstantTok{TRUE}\NormalTok{) }\SpecialCharTok{{-}} \FunctionTok{min}\NormalTok{(x, }\AttributeTok{na.rm =} \ConstantTok{TRUE}\NormalTok{)))}
\NormalTok{\}}

\CommentTok{\# Aplicar normalización a las variables numéricas}
\NormalTok{datos\_norm }\OtherTok{\textless{}{-}}\NormalTok{ variables\_numericas }\SpecialCharTok{\%\textgreater{}\%} \FunctionTok{mutate}\NormalTok{(}\FunctionTok{across}\NormalTok{(}\FunctionTok{everything}\NormalTok{(), normalizar))}

\CommentTok{\# Formato largo para los boxplots}
\NormalTok{datos\_norm\_long }\OtherTok{\textless{}{-}}\NormalTok{ datos\_norm }\SpecialCharTok{\%\textgreater{}\%} 
  \FunctionTok{pivot\_longer}\NormalTok{(}\AttributeTok{cols =} \FunctionTok{everything}\NormalTok{(), }\AttributeTok{names\_to =} \StringTok{"Variable"}\NormalTok{, }\AttributeTok{values\_to =} \StringTok{"Valor\_Normalizado"}\NormalTok{)}

\CommentTok{\# Graficar Boxplots unificados}
\NormalTok{p\_box }\OtherTok{\textless{}{-}} \FunctionTok{ggplot}\NormalTok{(datos\_norm\_long, }\FunctionTok{aes}\NormalTok{(}\AttributeTok{x =} \FunctionTok{reorder}\NormalTok{(Variable, Valor\_Normalizado, }\AttributeTok{FUN =}\NormalTok{ median, }\AttributeTok{na.rm =} \ConstantTok{TRUE}\NormalTok{), }\AttributeTok{y =}\NormalTok{ Valor\_Normalizado)) }\SpecialCharTok{+}
  \FunctionTok{geom\_boxplot}\NormalTok{(}\AttributeTok{fill =} \StringTok{"\#2ecc71"}\NormalTok{, }\AttributeTok{color =} \StringTok{"\#27ae60"}\NormalTok{, }\AttributeTok{outlier.color =} \StringTok{"\#e74c3c"}\NormalTok{, }\AttributeTok{outlier.alpha =} \FloatTok{0.5}\NormalTok{) }\SpecialCharTok{+}
  \FunctionTok{coord\_flip}\NormalTok{() }\SpecialCharTok{+} \CommentTok{\# Voltear para mejorar la lectura de las etiquetas}
  \FunctionTok{theme\_minimal}\NormalTok{() }\SpecialCharTok{+}
  \FunctionTok{labs}\NormalTok{(}
    \AttributeTok{title =} \StringTok{"Diagramas de Caja (Variables Normalizadas [0,1])"}\NormalTok{,}
    \AttributeTok{x =} \StringTok{"Variable Numérica"}\NormalTok{,}
    \AttributeTok{y =} \StringTok{"Valor Normalizado"}
\NormalTok{  )}

\CommentTok{\# Mostrar y guardar gráfico}
\FunctionTok{print}\NormalTok{(p\_box)}
\end{Highlighting}
\end{Shaded}

\begin{center}\includegraphics{Taller1_files/figure-latex/boxplots_normalizados-1} \end{center}

\begin{Shaded}
\begin{Highlighting}[]
\FunctionTok{ggsave}\NormalTok{(}\StringTok{"resultados taller 1/Boxplots.png"}\NormalTok{, }\AttributeTok{plot =}\NormalTok{ p\_box, }\AttributeTok{width =} \DecValTok{8}\NormalTok{, }\AttributeTok{height =} \DecValTok{10}\NormalTok{)}
\end{Highlighting}
\end{Shaded}

\begin{center}\rule{0.5\linewidth}{0.5pt}\end{center}

\subsection{5. Relaciones Bivariadas (Diagramas de
Dispersión)}\label{relaciones-bivariadas-diagramas-de-dispersiuxf3n}

Analizamos las interdependencias críticas entre indicadores
seleccionados, evaluando su trayectoria y posibles tendencias lineales.

\begin{Shaded}
\begin{Highlighting}[]
\CommentTok{\# Definir los 5 pares de variables solicitados}
\NormalTok{pares }\OtherTok{\textless{}{-}} \FunctionTok{list}\NormalTok{(}
  \FunctionTok{c}\NormalTok{(}\StringTok{"COBERTURA\_NETA"}\NormalTok{, }\StringTok{"COBERTURA\_BRUTA"}\NormalTok{),}
  \FunctionTok{c}\NormalTok{(}\StringTok{"APROBACIÓN\_SECUNDARIA"}\NormalTok{, }\StringTok{"REPROBACIÓN\_SECUNDARIA"}\NormalTok{),}
  \FunctionTok{c}\NormalTok{(}\StringTok{"DESERCIÓN"}\NormalTok{, }\StringTok{"DESERCIÓN\_SECUNDARIA"}\NormalTok{),}
  \FunctionTok{c}\NormalTok{(}\StringTok{"REPITENCIA"}\NormalTok{, }\StringTok{"REPITENCIA\_PRIMARIA"}\NormalTok{),}
  \FunctionTok{c}\NormalTok{(}\StringTok{"APROBACIÓN"}\NormalTok{, }\StringTok{"REPROBACIÓN"}\NormalTok{) }\CommentTok{\# Quinto par añadido para complementar el análisis}
\NormalTok{)}

\CommentTok{\# Generar gráficos secuenciales y exportarlos}
\ControlFlowTok{for}\NormalTok{ (i }\ControlFlowTok{in} \DecValTok{1}\SpecialCharTok{:}\FunctionTok{length}\NormalTok{(pares)) \{}
\NormalTok{  var1 }\OtherTok{\textless{}{-}}\NormalTok{ pares[[i]][}\DecValTok{1}\NormalTok{]}
\NormalTok{  var2 }\OtherTok{\textless{}{-}}\NormalTok{ pares[[i]][}\DecValTok{2}\NormalTok{]}
  
  \ControlFlowTok{if}\NormalTok{(var1 }\SpecialCharTok{\%in\%} \FunctionTok{names}\NormalTok{(datos) }\SpecialCharTok{\&}\NormalTok{ var2 }\SpecialCharTok{\%in\%} \FunctionTok{names}\NormalTok{(datos)) \{}
\NormalTok{    p }\OtherTok{\textless{}{-}} \FunctionTok{ggplot}\NormalTok{(datos, }\FunctionTok{aes\_string}\NormalTok{(}\AttributeTok{x =}\NormalTok{ var1, }\AttributeTok{y =}\NormalTok{ var2)) }\SpecialCharTok{+}
      \FunctionTok{geom\_point}\NormalTok{(}\AttributeTok{alpha =} \FloatTok{0.5}\NormalTok{, }\AttributeTok{color =} \StringTok{"\#2c3e50"}\NormalTok{, }\AttributeTok{size =} \DecValTok{2}\NormalTok{) }\SpecialCharTok{+}
      \FunctionTok{geom\_smooth}\NormalTok{(}\AttributeTok{method =} \StringTok{"lm"}\NormalTok{, }\AttributeTok{color =} \StringTok{"\#e74c3c"}\NormalTok{, }\AttributeTok{se =} \ConstantTok{FALSE}\NormalTok{, }\AttributeTok{linetype =} \StringTok{"dashed"}\NormalTok{) }\SpecialCharTok{+}
      \FunctionTok{labs}\NormalTok{(}
        \AttributeTok{title =} \FunctionTok{paste}\NormalTok{(}\StringTok{"Relación entre"}\NormalTok{, var1, }\StringTok{"y"}\NormalTok{, var2), }
        \AttributeTok{x =}\NormalTok{ var1, }
        \AttributeTok{y =}\NormalTok{ var2}
\NormalTok{      ) }\SpecialCharTok{+}
      \FunctionTok{theme\_minimal}\NormalTok{()}
    
    \FunctionTok{print}\NormalTok{(p)}
    \FunctionTok{ggsave}\NormalTok{(}\FunctionTok{paste0}\NormalTok{(}\StringTok{"resultados taller 1/Dispersion\_"}\NormalTok{, var1, }\StringTok{"\_vs\_"}\NormalTok{, var2, }\StringTok{".png"}\NormalTok{), }\AttributeTok{plot =}\NormalTok{ p, }\AttributeTok{width =} \DecValTok{6}\NormalTok{, }\AttributeTok{height =} \DecValTok{4}\NormalTok{)}
\NormalTok{  \}}
\NormalTok{\}}
\end{Highlighting}
\end{Shaded}

\begin{center}\includegraphics{Taller1_files/figure-latex/diagramas_dispersion-1} \end{center}

\begin{center}\includegraphics{Taller1_files/figure-latex/diagramas_dispersion-2} \end{center}

\begin{center}\includegraphics{Taller1_files/figure-latex/diagramas_dispersion-3} \end{center}

\begin{center}\includegraphics{Taller1_files/figure-latex/diagramas_dispersion-4} \end{center}

\begin{center}\includegraphics{Taller1_files/figure-latex/diagramas_dispersion-5} \end{center}

\begin{center}\rule{0.5\linewidth}{0.5pt}\end{center}

\subsection{6. Matriz de Correlación Completa
(Pearson)}\label{matriz-de-correlaciuxf3n-completa-pearson}

Asociación lineal cruzada entre todas las variables cuantitativas del
sistema educativo.

\begin{Shaded}
\begin{Highlighting}[]
\CommentTok{\# Matriz de Correlación ignorando valores faltantes}
\NormalTok{matriz\_corr }\OtherTok{\textless{}{-}} \FunctionTok{cor}\NormalTok{(variables\_numericas, }\AttributeTok{use =} \StringTok{"complete.obs"}\NormalTok{)}

\CommentTok{\# Guardar matriz de correlación en CSV}
\FunctionTok{write.csv}\NormalTok{(matriz\_corr, }\StringTok{"resultados taller 1/Matriz\_Correlacion.csv"}\NormalTok{)}

\CommentTok{\# Transformación de la matriz para el mapa de calor (Heatmap)}
\NormalTok{melted\_corr }\OtherTok{\textless{}{-}} \FunctionTok{melt}\NormalTok{(matriz\_corr)}

\CommentTok{\# Graficar el Heatmap}
\NormalTok{p\_corr }\OtherTok{\textless{}{-}} \FunctionTok{ggplot}\NormalTok{(melted\_corr, }\FunctionTok{aes}\NormalTok{(Var1, Var2, }\AttributeTok{fill =}\NormalTok{ value)) }\SpecialCharTok{+}
  \FunctionTok{geom\_tile}\NormalTok{(}\AttributeTok{color =} \StringTok{"white"}\NormalTok{) }\SpecialCharTok{+}
  \FunctionTok{scale\_fill\_gradient2}\NormalTok{(}
    \AttributeTok{low =} \StringTok{"\#2980b9"}\NormalTok{, }\AttributeTok{high =} \StringTok{"\#c0392b"}\NormalTok{, }\AttributeTok{mid =} \StringTok{"white"}\NormalTok{, }
    \AttributeTok{midpoint =} \DecValTok{0}\NormalTok{, }\AttributeTok{limit =} \FunctionTok{c}\NormalTok{(}\SpecialCharTok{{-}}\DecValTok{1}\NormalTok{,}\DecValTok{1}\NormalTok{), }\AttributeTok{name =} \StringTok{"Pearson (r)"}
\NormalTok{  ) }\SpecialCharTok{+}
  \FunctionTok{theme\_minimal}\NormalTok{() }\SpecialCharTok{+}
  \FunctionTok{theme}\NormalTok{(}
    \AttributeTok{axis.text.x =} \FunctionTok{element\_text}\NormalTok{(}\AttributeTok{angle =} \DecValTok{45}\NormalTok{, }\AttributeTok{vjust =} \DecValTok{1}\NormalTok{, }\AttributeTok{hjust =} \DecValTok{1}\NormalTok{, }\AttributeTok{size =} \DecValTok{8}\NormalTok{),}
    \AttributeTok{axis.text.y =} \FunctionTok{element\_text}\NormalTok{(}\AttributeTok{size =} \DecValTok{8}\NormalTok{),}
    \AttributeTok{axis.title =} \FunctionTok{element\_blank}\NormalTok{(),}
    \AttributeTok{panel.grid.major =} \FunctionTok{element\_blank}\NormalTok{()}
\NormalTok{  ) }\SpecialCharTok{+}
  \FunctionTok{labs}\NormalTok{(}\AttributeTok{title =} \StringTok{"Matriz de Correlación (Todas las Variables Numéricas)"}\NormalTok{)}

\CommentTok{\# Mostrar y guardar gráfico}
\FunctionTok{print}\NormalTok{(p\_corr)}
\end{Highlighting}
\end{Shaded}

\begin{center}\includegraphics{Taller1_files/figure-latex/matriz_correlacion-1} \end{center}

\begin{Shaded}
\begin{Highlighting}[]
\FunctionTok{ggsave}\NormalTok{(}\StringTok{"resultados taller 1/Heatmap\_Correlacion.png"}\NormalTok{, }\AttributeTok{plot =}\NormalTok{ p\_corr, }\AttributeTok{width =} \DecValTok{10}\NormalTok{, }\AttributeTok{height =} \DecValTok{10}\NormalTok{)}
\end{Highlighting}
\end{Shaded}

\begin{center}\rule{0.5\linewidth}{0.5pt}\end{center}

\subsection{7. Conclusiones del
Análisis}\label{conclusiones-del-anuxe1lisis}

Con base en los gráficos y el modelo estadístico construido, se destaca
lo siguiente:

\begin{enumerate}
\def\labelenumi{\arabic{enumi}.}
\tightlist
\item
  \textbf{Variables Geográficas vs.~Métricas:} Se aplicó la separación
  dictada por el diccionario de datos: los Códigos DANE municipales y
  departamentales fueron tratados como dimensiones categóricas, dejando
  un espacio analítico puramente enfocado en tasas, porcentajes y
  coberturas.
\item
  \textbf{Distribuciones (Histogramas):} Se observa que las variables de
  \emph{Aprobación} tienen un marcado sesgo hacia la derecha (alta
  concentración cerca al 100\%), mientras que las métricas de
  \emph{Deserción} y \emph{Reprobación} tienen distribuciones sesgadas
  fuertemente hacia la izquierda, denotando valores bajos mayoritarios
  con una larga cola de casos críticos.
\item
  \textbf{Identificación de Atípicos (Boxplots):} En el plano
  normalizado, la \emph{Deserción} y la \emph{Repitencia} presentan una
  gran cantidad de municipios por fuera de los ``bigotes'' superiores,
  indicando enclaves territoriales con crisis educativas severas.
\item
  \textbf{Relaciones Lineales Bivariadas:}

  \begin{itemize}
  \tightlist
  \item
    Se confirma una correlación negativa perfecta entre Aprobación y
    Reprobación, como es naturalmente esperado.
  \item
    La Deserción Secundaria tracciona sustancialmente el indicador de
    Deserción General, indicando que el abandono escolar es prevalente
    en la educación media/secundaria más que en preescolar o primaria.
  \end{itemize}
\item
  \textbf{Matriz de Correlación Multivariada:} Se detectan bloques
  altamente interconectados. Las variables de Aprobación por niveles
  (Transición, Primaria, Secundaria) correlacionan positivamente entre
  sí, y negativamente contra cualquier tipo de Reprobación y Deserción,
  con un patrón cromático en la matriz que divide visualmente las
  métricas positivas frente a los déficits del sistema educativo.
\end{enumerate}

\end{document}
